\chapter{Hashing de contenido y direccionamiento}
\label{cap:hashing}

El TypeGraph del cap.~\ref{cap:typegraph} fija qué tipos viven en el
grafo. Este capítulo agrega la \emph{identidad por contenido}: cada
nodo tiene un identificador derivado del hash criptográfico de su
contenido serializado. Esa decisión apenas perceptible en la API tiene
consecuencias profundas: deduplicación intrínseca, snapshots gratis,
verificación incremental, sharing entre dominios sin bolt-on.

\section{El hash de un nodo}

\begin{defn}[Content hash]
\label{def:content-hash}
Para todo nodo $v \in V$, su \emph{content hash} es
\[
  \hashof{v} \;=\; \blake\bigl(\tau_V(v) \,\Vert\, \mathrm{serialize}(v)\bigr)
\]
donde $\Vert$ denota concatenación, $\tau_V(v)$ es un byte-tag que
identifica el tipo del nodo (8-bit) y $\mathrm{serialize}$ es la
serialización canónica del contenido.
\end{defn}

La canonicidad de $\mathrm{serialize}$ es crítica: dos nodos con el
mismo contenido lógico pero distinta serialización se hashean a valores
distintos, lo que rompería la deduplicación. \citaCrate{sotfs-graph}
usa una serialización determinista (campos en orden fijo, sin padding,
sin reordering) basada en \texttt{serde\_json} configurado para emitir
keys ordenadas. Para los \NBlk{} en particular, el contenido es la
sucesión cruda de bytes, sin metadatos: dos archivos con el mismo
contenido producen el mismo \texttt{BlockId}.

\paragraph{Algoritmo.}

\begin{algorithm}[H]
\SetAlgoLined
\KwIn{Nodo $v$ con tipo $\tau_V(v)$ y contenido $\mathrm{content}(v)$}
\KwOut{Hash $h \in \{0,1\}^{256}$}
\caption{\textsc{NodeHash}}\label{alg:node-hash}

$\mathrm{hasher} \gets \blake.\textsc{New}()$\;
$\mathrm{hasher}.\textsc{Update}(\tau_V(v).\textsc{ToBytes}())$\;
$\mathrm{hasher}.\textsc{Update}(\mathrm{Serialize}(\mathrm{content}(v)))$\;
\Return{$\mathrm{hasher}.\textsc{Finalize}()$}\;
\end{algorithm}

\section{Identidad por contenido y deduplicación}

\begin{invariante}[Identidad por contenido]
\label{inv:identidad-contenido}
Dos nodos con el mismo $\hashof{\cdot}$ son el mismo nodo en el
TypeGraph: el filesystem mantiene una sola copia. Identidad estructural
implica identidad de almacenamiento.
\end{invariante}

La consecuencia operativa: dos archivos cuya escritura produce los
mismos bloques no duplican espacio. La operación \opwrite{}
(cap.~\ref{cap:dpo}) computa el \texttt{BlockId} de cada bloque a
escribir; si ya existe en \citaCrate{sotfs-storage}, simplemente añade
una arista \epointsto{} y bumpea el \texttt{refcount} del bloque
existente.

\begin{figure}[h]
\centering
\begin{tikzpicture}[node distance=10mm and 16mm]
  \node[nino] (i1) {Inode A\\\texttt{file1.txt}};
  \node[nino, right=30mm of i1] (i2) {Inode B\\\texttt{file2.txt}};
  \node[nblk, below=14mm of i1, xshift=8mm] (b1) {Block\\\hashof{·}=A1};
  \node[nblk, below=14mm of i1, xshift=30mm] (b2) {Block compartido\\\hashof{·}=B2};
  \node[nblk, below=14mm of i2, xshift=12mm] (b3) {Block\\\hashof{·}=C3};

  \draw[arr-points] (i1) -- (b1);
  \draw[arr-points] (i1) -- (b2);
  \draw[arr-points] (i2) -- (b2);
  \draw[arr-points] (i2) -- (b3);
\end{tikzpicture}
\caption{Sharing intrínseco: \texttt{file1.txt} y \texttt{file2.txt}
comparten el bloque \texttt{B2} por construcción porque ambos lo
producen con el mismo hash. No es un hard link; es la semántica natural
del Merkle DAG. El \texttt{refcount} de \texttt{B2} es 2.}
\label{fig:hashing-sharing}
\end{figure}

\section{Verificación incremental}

\begin{teorema}[Verificación cripto-derivada]
\label{teo:verificacion-incremental}
Dado el hash de la raíz de un sub-DAG $G_R$, el contenido de cada nodo
del sub-DAG se puede verificar sin material adicional, en tiempo $O(N
\cdot \mathrm{cost}(\blake))$, donde $N$ es la cantidad de nodos en el
sub-DAG.
\end{teorema}

\begin{proof}[Esquema]
Por inducción sobre la profundidad del sub-DAG. La raíz se verifica
recomputando $\hashof{r}$ y comparando con el dado. Cada hijo aparece
referenciado por su hash dentro del contenido del padre; recomputar
$\hashof{\mathrm{hijo}}$ y comparar con el referenciado verifica el
hijo. La consistencia entre padre y hijo es automática porque la
referencia (hash) está en el contenido del padre.
\end{proof}

Esa propiedad es la base de los snapshots: un snapshot es \emph{nada
más que} un puntero al hash de un nodo \NVer{}. Para verificar que el
snapshot no fue manipulado, se recomputa el hash del nodo \NVer{} y se
compara.

\section{Snapshots gratis}

Bajo el modelo content-addressed, hacer un snapshot de un FS de 1\,TiB
con \(10^6\) archivos cuesta exactamente \emph{32 bytes}: el hash del
nodo \NVer{} que apunta a la raíz. Comparar con \texttt{ext4} con
LVM-snapshots, donde el costo es proporcional al número de bloques
modificados desde el snapshot original (vía COW): $O(\Delta)$ bytes y
$O(\Delta)$ overhead de tracking.

La desventaja: garbage collection. Cada snapshot ``pin'' a un sub-DAG;
nodos no alcanzables desde ningún snapshot son candidatos a GC. La
implementación es mark-and-sweep desde los roots de snapshot;
cap.~\ref{cap:persistencia} desarrolla los detalles.

\begin{trampa}[BLAKE3 collision (no observada en práctica)]
\label{tr:blake3-collision}
La \invref{inv:identidad-contenido} asume que los hashes son únicos.
\blake{} con 256 bits tiene resistencia a colisión de $2^{128}$
(birthday bound); la probabilidad de un ciclo accidental por colisión
es \emph{efectivamente cero}. Pero un atacante con poder de construir
\emph{colisiones intencionales} (problema abierto en \blake{}) podría
engañar al filesystem para que un nodo se auto-referencie a través de
su contenido.

\textbf{Fix preventivo.} Chequeo runtime que un nodo recién insertado
no contiene su propio hash entre sus referencias. Costo $O(|\text{children}|)$
por inserción. Reductible si hace falta.

\textbf{Lección.} Las propiedades cripto-derivadas son fuertes pero no
infalibles. Un check runtime trivial (anti-self-reference) preserva la
invariante incluso bajo collision attack.
\end{trampa}

\section{Sharing entre dominios}

\begin{defn}[Dominio]
\label{def:dominio}
Un \emph{dominio} es una partición lógica del namespace del filesystem,
identificada por un \texttt{domain\_id} (un \texttt{u64}). Los caps
están parametrizados por el \texttt{domain\_id} del operador
(\citaLine{sotfs-graph/src/types.rs}{189}).
\end{defn}

Dos dominios distintos pueden referenciar el mismo \NBlk{} porque el
\NBlk{} se nombra por hash; no hay copia ni protocolo cross-domain. La
propiedad emerge del modelo. La consecuencia: dos contenedores
distintos en sotFS que comparten la misma imagen base pueden compartir
bloques sin sobrecosto.

\section{Comparación con árbol POSIX}

\begin{table}[h]
\centering
\renewcommand{\arraystretch}{1.2}
\begin{tabularx}{\linewidth}{l X X}
  \toprule
   & \textbf{Árbol POSIX (\texttt{ext4})} & \textbf{Merkle DAG (sotFS)} \\
  \midrule
  Identidad   & \texttt{InodeId} arbitrario, allocator-dependent & $\hashof{\cdot}$ del contenido \\
  Dedup       & manual / scrub & intrínseca \\
  Snapshot    & LVM/COW, $O(\Delta)$ bytes & 32 bytes \\
  Sharing     & hard link, simulink, reflink (extensión) & emerge del modelo \\
  Verificación & no nativa & $O(N \cdot \mathrm{cost}(\blake))$ recursiva \\
  Mutación    & in-place a bloque & inmutable; nuevo nodo, padre actualizado \\
  GC          & no necesario (sin sharing) & necesario; mark-and-sweep \\
  \bottomrule
\end{tabularx}
\caption{Tradeoffs de árbol POSIX vs.\ Merkle DAG.}
\label{tab:hash-tradeoffs}
\end{table}

\section{Inmutabilidad de nodo}

\begin{invariante}[Inmutabilidad estructural]
\label{inv:inmutabilidad}
Un nodo nunca se modifica \emph{in situ}. Una mutación lógica
(\opwrite, \texttt{truncate}, \texttt{chmod}) crea un nodo nuevo con el
contenido nuevo y actualiza el padre. El nodo viejo sobrevive si está
referenciado por algún snapshot; si no, queda candidato a GC.
\end{invariante}

\section*{Resumen del capítulo}

\begin{itemize}
  \item Cada nodo del TypeGraph tiene un content hash $\hashof{v} =
    \blake(\tau_V(v) \Vert \mathrm{serialize}(v))$, calculado con
    serialización canónica.
  \item Identidad por contenido implica deduplicación intrínseca y
    sharing entre dominios sin bolt-on.
  \item Verificación incremental: dado el hash raíz, todo el sub-DAG
    se verifica en $O(N \cdot \mathrm{cost}(\blake))$.
  \item Snapshots gratis: 32 bytes para snapshot un FS de 1\,TiB.
  \item Inmutabilidad estructural: las mutaciones crean nodos nuevos,
    no modifican in-place.
  \item Trampa: colisiones intencionales de \blake{} son un problema
    abierto; chequeo anti-self-reference como mitigación barata.
\end{itemize}

\section*{Ejercicios}

\begin{ejercicio}[\dificultad{$\star$}]
\label{ej:hash-1}
¿Cuál es la longitud en bytes del hash \blake{} usado por sotFS?
¿Cuál es la cota de colisión por birthday bound?
\end{ejercicio}

\begin{ejercicio}[\dificultad{$\star$}]
\label{ej:hash-2}
Indique tres consecuencias del modelo content-addressed sobre la
operación \opwrite{} en sotFS.
\end{ejercicio}

\begin{ejercicio}[\dificultad{$\star\star$}]
\label{ej:hash-3}
Dos archivos de 1\,MiB tienen 80\% de contenido en común, dividido en
bloques de 4\,KiB. Calcule el espacio total ocupado en sotFS y compare
con \texttt{ext4} sin reflinks.
\end{ejercicio}

\begin{ejercicio}[\dificultad{$\star\star$}]
\label{ej:hash-4}
Argumente por qué la canonicidad de la serialización es crítica.
Construya un escenario donde dos serializaciones distintas del mismo
contenido lógico romperían la deduplicación.
\end{ejercicio}

\begin{ejercicio}[\dificultad{$\star\star$}]
\label{ej:hash-5}
Diseñe un mecanismo de \emph{garbage collection} para sotFS basado en
mark-and-sweep desde los snapshot roots. ¿Cómo manejás writes
concurrentes que crean nodos durante el mark? Pista: tri-color marking
+ write barrier sobre roots.
\end{ejercicio}
